%!TEX root = tfg.tex
\chapter{Plan de Dirección del Proyecto}

\begin{quotation}[General and President]{Dwight D. Eisenhower (1890--1969)}
Plans are useless, but planning is indispensable.
\end{quotation}

\begin{abstract}
En este capítulo se define el Plan de Dirección del Proyecto de \textit{Via Inferni},
tomando como referencia el marco PMBOK y seleccionando los procesos que resultan útiles
para un Trabajo de Fin de Grado desarrollado por dos estudiantes de Ingeniería Informática.
El propósito del capítulo es establecer el marco de dirección que integra los planes
subsidiarios, las líneas base, los mecanismos de control y los criterios de cierre del
proyecto \cite{pmbok2021}.
\end{abstract}

\section{Propósito del Plan de Dirección}
El Plan de Dirección del Proyecto actúa como documento integrador de la planificación. De
acuerdo con PMBOK \cite{pmbok2021}, su función es coordinar los planes de gestión, las
líneas base y los procedimientos de seguimiento que permiten comparar el trabajo previsto
con el trabajo realmente ejecutado.

En el contexto de \textit{Via Inferni}, el plan se aplica de forma proporcional. El
proyecto no requiere una estructura documental pesada, pero sí necesita criterios claros
para controlar el alcance, ordenar el cronograma, registrar decisiones técnicas y justificar
las desviaciones que aparezcan durante el desarrollo. Por ello, este plan se centra en los
elementos que aportan trazabilidad y capacidad de decisión al equipo.

La planificación persigue tres objetivos principales:

\begin{itemize}
    \item delimitar un alcance viable para el videojuego dentro del tiempo disponible;
    \item organizar el desarrollo mediante incrementos jugables, verificables y documentados;
    \item mantener coherencia entre la planificación, el código fuente, la documentación y
    los resultados obtenidos en cada iteración.
\end{itemize}

\section{Enfoque de Dirección basado en PMBOK}
El proyecto adopta PMBOK como marco de referencia y selecciona las áreas que aportan valor
real a un equipo reducido y a un entorno académico. Se desarrollan como planes subsidiarios
aquellas áreas que tienen impacto directo en la ejecución del TFG: alcance, cronograma,
comunicaciones, riesgos, calidad y costes. Otras áreas, como interesados, recursos y
adquisiciones, se tratan de forma integrada o se justifican como no aplicables en el
capítulo de selecciones y excepciones.

La gestión de la integración se materializa en este propio Plan de Dirección. Su función es
asegurar que las decisiones sobre alcance, calendario, calidad, riesgos, configuración y
cierre se tomen de forma coherente. En consecuencia, cualquier cambio que afecte a una
línea base debe analizarse considerando su impacto sobre el conjunto del proyecto, y no solo
sobre la tarea concreta en la que aparece.

\subsection{Planes subsidiarios incluidos}
Los planes de gestión que forman parte de este Plan de Dirección son los siguientes:

\begin{itemize}
    \item \textbf{Plan de gestión del alcance:} define cómo se delimita el producto y cómo
    se aceptan, posponen o descartan nuevas funcionalidades.
    \item \textbf{Plan de gestión del cronograma:} organiza el trabajo en fases e
    iteraciones, fija hitos y establece la estructura interna de cada ciclo de desarrollo.
    \item \textbf{Plan de gestión de comunicaciones:} establece los canales, reuniones y
    mecanismos de registro de decisiones.
    \item \textbf{Plan de gestión de riesgos:} identifica las principales incertidumbres del
    proyecto y las respuestas previstas.
    \item \textbf{Plan de gestión de la calidad:} fija los criterios que debe cumplir una
    funcionalidad para considerarse terminada.
    \item \textbf{Plan de gestión de costes:} estima el valor económico del esfuerzo y de los
    recursos empleados.
\end{itemize}

El desarrollo detallado de estos planes se presenta en el capítulo siguiente.

\subsection{Líneas base del proyecto}
El proyecto queda apoyado por tres líneas base principales:

\begin{itemize}
    \item \textbf{Línea base del alcance:} formada por el modelo general del producto y la
    lista de funcionalidades priorizadas.
    \item \textbf{Línea base del cronograma:} formada por las fases del proyecto, las ocho
    iteraciones de desarrollo y los hitos de seguimiento y cierre.
    \item \textbf{Línea base de costes:} formada por la estimación de horas, coste de
    personal, amortización de equipamiento y costes indirectos.
\end{itemize}

Estas líneas base permiten evaluar desviaciones, justificar replanificaciones y mantener
una referencia común durante la ejecución.

\section{Ciclo de vida del proyecto}
El ciclo de vida del proyecto se organiza en tres fases:

\begin{itemize}
    \item \textbf{Planificación:} definición de objetivos, alcance, metodología,
    herramientas, repositorio, backlog inicial y planes de gestión.
    \item \textbf{Ejecución, seguimiento y control:} desarrollo iterativo del videojuego,
    validación de funcionalidades, revisión de riesgos, control de calidad y actualización
    de la documentación.
    \item \textbf{Cierre:} consolidación del producto, revisión final de la memoria,
    preparación de manuales y formalización de conclusiones.
\end{itemize}

Esta estructura combina una planificación inicial suficiente con una ejecución iterativa que
permite ajustar el trabajo conforme se obtiene conocimiento técnico del videojuego.

\section{Marco metodológico de desarrollo}
La metodología de desarrollo elegida para organizar el trabajo técnico es
\textit{Feature Driven Development} (FDD) \cite{palmer2002fdd}. Esta metodología se utiliza
como soporte operativo del plan de alcance y del plan de cronograma, no como sustituto del
marco de dirección PMBOK.

FDD resulta adecuada para el proyecto porque permite descomponer el videojuego en
funcionalidades pequeñas, priorizables y verificables. En lugar de abordar el producto como
un bloque único, el desarrollo se articula mediante incrementos que pueden analizarse,
diseñarse, implementarse, probarse e integrarse en iteraciones sucesivas.

El proceso general se resume en cinco pasos: desarrollo de un modelo general, construcción
de la lista de funcionalidades, planificación por funcionalidad, diseño por funcionalidad y
construcción por funcionalidad. Este flujo se representa en la Figura \ref{fig:fdd_flow}. La
aplicación concreta de estos pasos al alcance y al cronograma se desarrolla en los planes de
gestión correspondientes.

\begin{figure}[htb]
    \centering
    \begin{tikzpicture}[node distance=1.2cm, auto,
        block/.style={rectangle, draw, fill=gray!5, text width=6cm, text centered, rounded corners, minimum height=1cm, font=\small},
        line/.style={draw, -latex, thick}]
        \node [block] (step1) {\textbf{1. Desarrollo de un modelo general} \\ (agrupación de dominios funcionales)};
        \node [block, below=of step1] (step2) {\textbf{2. Construcción de la lista de funcionalidades} \\ (desglose del backlog)};
        \node [block, below=of step2] (step3) {\textbf{3. Planificación por funcionalidad} \\ (priorización e hitos)};
        \node [block, below=of step3, fill=blue!5] (step4) {\textbf{4. Diseño por funcionalidad} \\ (contratos y lógica técnica)};
        \node [block, below=of step4, fill=blue!5] (step5) {\textbf{5. Construcción por funcionalidad} \\ (codificación e integración)};

        \path [line] (step1) -- (step2);
        \path [line] (step2) -- (step3);
        \path [line] (step3) -- (step4);
        \path [line] (step4) -- (step5);
        \path [line] (step5.east) -- ++(1,0) -- ++(0,2.2) -- (step4.east);
        \draw[dashed, blue!40, thick] ([xshift=-0.5cm, yshift=0.3cm]step4.north west) rectangle ([xshift=2.5cm, yshift=-0.3cm]step5.south east);
        \node[blue!60, rotate=90, anchor=south] at ([xshift=2.5cm]step5.east) {Fase iterativa};
    \end{tikzpicture}
    \caption{Flujo de procesos de la metodología FDD aplicado al desarrollo del sistema.}
    \label{fig:fdd_flow}
\end{figure}

\section{Gestión de la configuración, código fuente y cambios}
La gestión de la configuración agrupa el control del repositorio, la organización del código
fuente y el procedimiento de gestión de cambios. Estos elementos se tratan de forma conjunta
porque comparten un mismo objetivo: mantener la integridad del producto y la trazabilidad de
las decisiones durante todo el desarrollo.

El repositorio se gestiona mediante una estrategia híbrida que combina la organización
estructural de \textit{GitFlow} \cite{driessen2010gitflow} con la integración frecuente
propia de \textit{Trunk-Based Development} \cite{trunkbasedIntro}. La rama \texttt{main}
representa el estado estable del proyecto, mientras que \texttt{develop} actúa como punto de
integración del trabajo validado. Las ramas de funcionalidad son temporales y se crean para
resolver tareas concretas del backlog.

Como criterio de trazabilidad se adopta el estándar \textit{Conventional Commits}
\cite{conventionalCommits10}. Este criterio permite identificar la intención de cada cambio,
distinguir funcionalidades, correcciones y ajustes documentales, y relacionar el historial
del repositorio con el avance de las iteraciones.

El control de cambios se aplica de forma proporcional a su impacto:

\begin{itemize}
    \item los ajustes menores que no afectan al alcance, al cronograma ni a la calidad se
    resuelven dentro de la iteración correspondiente;
    \item los cambios que modifican una funcionalidad planificada, desplazan hitos o generan
    nuevos riesgos deben acordarse entre los dos miembros del equipo;
    \item las decisiones relevantes se registran en el backlog, en GitHub Projects o en la
    documentación de la iteración, según corresponda.
\end{itemize}

Este procedimiento evita incorporar trabajo no planificado sin evaluar su impacto. La
decisión de aceptar, dividir, posponer o descartar una funcionalidad se toma comparando su
valor para el núcleo jugable con el esfuerzo técnico y temporal que exige.

\section{Herramientas y tecnologías}
El proyecto se apoya en un conjunto reducido de herramientas técnicas y de gestión. Su
selección responde a criterios de adecuación al tipo de videojuego, madurez técnica del
equipo, trazabilidad, coste y facilidad de integración con la memoria.

Para el motor de desarrollo se valoraron alternativas como Godot, Unreal Engine y Unity.
Godot ofrecía una solución ligera y de código abierto, adecuada para proyectos 2D, pero el
equipo consideró que Unity podía ofrecer más posibilidades de desarrollo y una transición
más natural gracias a la experiencia previa con C\#. Unreal Engine proporcionaba una
capacidad técnica elevada, especialmente en proyectos 3D, pero su complejidad y orientación
principal hacían menos eficiente su uso para un \textit{roguelike} 2D de alcance académico.
Finalmente, se seleccionó Unity por su equilibrio entre capacidad técnica, documentación,
curva de aprendizaje asumible, gestión de escenas, soporte para interfaz y velocidad de
prototipado \cite{unityLearnGettingStarted}.

Para el control de versiones se consideraron flujos basados en repositorios locales,
GitLab, Bitbucket y GitHub. Se eligieron Git y GitHub porque eran herramientas ya conocidas
por el equipo, ofrecían integración directa entre repositorio, incidencias, revisión de
cambios y tablero Kanban, y facilitaban la trazabilidad del desarrollo en la memoria
\cite{githubIssuesProjects}. GitHub Projects se utiliza como tablero de seguimiento por
estar integrado con las tareas y los cambios del repositorio.

Para la comunicación se eligió Discord frente a alternativas como Microsoft Teams, Slack o
correo electrónico. La decisión se basó en su ligereza, disponibilidad para ambos miembros
del equipo y utilidad para resolver bloqueos operativos de forma rápida. Para el registro de
horas se seleccionó Clockify frente a una hoja de cálculo manual porque permite registrar
dedicación por tarea de forma más ordenada sin introducir una carga de gestión elevada
\cite{clockifyOfficial}.

Por último, la memoria se mantiene en LaTeX frente a procesadores de texto convencionales
porque el documento requiere modularidad, referencias cruzadas, tablas, figuras,
bibliografía y control estructural de un texto técnico extenso. Esta elección también
permite tratar la memoria como un artefacto versionado junto con el resto del proyecto
\cite{latexProject}.

\section{Criterios de seguimiento y cierre}
El seguimiento del proyecto se realiza de forma iterativa. Al cierre de cada ciclo se revisa
el estado de las funcionalidades, se comprueba el cumplimiento de la
\textit{Definition of Done}, se actualiza el backlog, se revisan riesgos y se documentan las
desviaciones relevantes.

El cierre del proyecto se considera alcanzado cuando el producto dispone de un núcleo
jugable verificable, la memoria refleja de forma coherente la evolución del desarrollo y los
materiales finales permiten instalar, comprender y evaluar el trabajo realizado.
